\section{Polynomial product norm suffices after logarithmic added depth}
\label{sec:product-targets}

The posterior representation ceiling has a partial converse for a
specified finite-width target representation. Polynomial product norm
alone does not balance the physical tanh layers, but small-prefix
linearization and unequal clone counts balance all retained hidden
layers. Only a polynomial readout cost then needs to be transmitted.

\begin{theorem}[Deep targets with polynomial product norm]
\label{thm:product-target-log-padding}
Suppose a width-$C$, depth-$D_0\ge2$ tanh network has global signed
margin at least one for $f$, and product of layer Frobenius norms at
most $P$. Put $P_0=\max\{1,P\}$ and
\[
 \tau^2=\frac1{16P_0D_0(d+1)^{3/2}},\qquad
 T=\frac{32CP_0^2}{\tau^2}.
\]
For every fixed $\zeta>0$, the original unrestricted Gaussian posterior
has sample complexity polynomial in
\[
 d,C,D_0,P_0,L,\varepsilon^{-1},\log(2/\delta)
\]
at every learner depth
\[
 L\ge D_0+\lceil\log T\rceil,
\]
simultaneously for every $N\ge C_L(f)$ and every
$0<\sigma^2\le N^{-(1+\zeta)}$. The same polynomial conclusion holds
under the alternative depth condition $L\ge2D_0$.
There is no bound on the individual weights of the specified target
representation and no internal robustness restriction.
\end{theorem}

\begin{proof}
Write $M_j$ for the actual Frobenius norm of its $j$th matrix, including
the readout, and $S_j=\prod_{i\le j}M_i$ for the prefix products.
All norms are positive, since the global margin excludes a zero map.
The augmented input norm is $R_0=\sqrt{d+1}$.
At every hidden layer, boundedness of tanh and the suffix Lipschitz
bound give
\[
 1\le |z(x)|\le\sqrt C\,\frac P{S_j},
 \qquad S_j\le P\sqrt C.
\]
Using an upper bound $P$ instead of the exact full product only weakens
this inequality.

Replace the tanh activation by the identity at each hidden layer whose
original prefix product satisfies $S_j<\tau$.
The inequality $|u-\tanh u|\le |u|^3/3$ and a telescoping argument give
total output change at most
\[
 \sum_{j:S_j<\tau}\frac{P}{S_j}
      \frac{(S_jR_0)^3}{3}
 \le\frac{PD_0\tau^2R_0^3}{3}\le\frac1{48}.
\]
The same prefix and suffix bounds apply throughout this telescoping,
since both identity and tanh are zero-preserving and one-Lipschitz.
Collapse the resulting consecutive linear matrices into their next
nonlinear layer or the final readout, then double the readout. The
result still computes $f$ with unit margin.

If every hidden activation was removed, $f$ is a Boolean halfspace,
and Theorem~\ref{thm:halfspace-unrestricted-all-depths} proves the claim.
Otherwise let $H\ge1$ be the number of retained hidden layers, with
original indices $j_1<\cdots<j_H$. Their retained prefix bounds satisfy
\[
 \tau\le S_{j_s}\le P_0\sqrt C.
\]
The first collapsed matrix has Frobenius norm at most $S_{j_1}$;
the $s$th has norm at most $S_{j_s}/S_{j_{s-1}}$; and the doubled
readout has norm at most $2P_0/S_{j_H}$. These are assigned norm bounds;
matrix multiplication can make the actual norms smaller without
affecting any subsequent upper bound. In particular every collapsed
physical weight is at most
\[
 R=\max\{1,2P_0\sqrt C/\tau\}.
\]

Define
\[
 a=\frac{d\tau^2}{8C},\qquad
 r_s=\frac{a}{dS_{j_s}^2}=\frac{\tau^2}{8C S_{j_s}^2},
 \qquad r_{\min}=\frac{\tau^2}{8C^2P_0^2}.
\]
Then $r_{\min}\le r_s\le1/(8C)$ and $a<1/8$.
At common width $N$, give every original neuron in retained layer $s$
exactly $c_s=\lceil Nr_s\rceil$ clones, divide outgoing weights by
$c_s$, and fill the remaining positions with zero neurons.
For
\[
 N\ge M:=\left\lceil\max\{2C,8D_0/r_{\min},2R,4T,4\}\right\rceil,
\]
the clone counts fit in width $N$ and
$Nr_s\le c_s\le Nr_s(1+1/(8D_0))$.
This preserves the collapsed network exactly. Its pre-readout energy
and readout energy obey
\[
 E_{\rm pre}\le a\left(1+\frac1{8D_0}\right)
 +(H-1)\left(1+\frac1{8D_0}\right)\le H,
\]
\[
 c^2=N\|v_N\|_2^2
 \le\frac{dS_{j_H}^2}{a}\frac{4P_0^2}{S_{j_H}^2}
 =\frac{32CP_0^2}{\tau^2}=T.
\]
For the first inequality, $H\le D_0-1$ and $a<1/8$ leave more than
enough slack to absorb all clone rounding.

Replace the readout by the scalar-transmission construction with
$m=L-H$ weighted layers, and double its final readout as in the
shallow-target proof. The resulting unit-margin witness has energy
\[
 K\le H+(m+15)T^{1/m}.
\]
The stated depth conditions ensure $m\ge2$.
If $L\ge D_0+\lceil\log T\rceil$, then $m\ge\log T$, and
$e^u\le1+2u$ for $0\le u\le1$ gives
\[
 K\le L+15+32\log T.
\]
Hence the Gaussian-distance parameter is bounded by
\[
 Q_*^2=((K+2)/L)^L\le e^{17}T^{32}.
\]
Under the alternative condition $L\ge2D_0$, one has $L/m\le2$ and
the direct estimate
\[
 K+2\le(L+17)T^{1/m},\qquad Q_*^2\le e^{17}T^2.
\]
In either case, $K,Q_*$ and $M$ are polynomial in the stated target
parameters and learner depth. The Gaussian-distance theorem therefore
gives the required sample bound for all $N\ge M$.

We check the actual smaller-width convention. The collapsed weights
are at most $R\le M/2$, and cloning can only reduce internal and
readout weights. The first scalar-transmission matrix has weights at
most $1.1R<M$. If $q^2=\max\{c^2,1\}^{1/m}\le T$, the remaining
scalar weights have magnitudes at most $2.2q/N$ and $2q/N$.
Since $M\ge4T\ge4$, these are less than $M$.
Thus the constructed width-$M$ target is admissible and $C_L(f)\le M$.
For $C_L(f)\le N<M$, clone a defining minimum-width target and double
its readout; its width and weights are bounded by a polynomial in $M$.
The finite-width theorem supplies the same polynomial sample form.
Intersect the two sample events, with failure probabilities $\delta/2$
each, to conclude.
\end{proof}

The theorem is polynomial in $P$ itself. For a generic bounded-weight
deep circuit, $P$ may still grow exponentially with its depth. Thus this
does not settle the unrestricted general small-circuit conjecture.
It does show that a polynomial product-norm representation is a
sufficient target condition after a logarithmic number of additional
layers, without imposing that condition on the support of the prior.
